\begin{table*}[!t]
\centering
\caption{\textbf{Provisional Three-Tier Interpretive Framework for LEN Resistance Surveillance.}}
\label{tab:framework}
\footnotesize
\begin{tabular}{p{1.2cm}p{2.8cm}p{4.2cm}p{3.8cm}}
\toprule
\textbf{Tier} & \textbf{Mutation / Pattern} & \textbf{Evidence Grade} & \textbf{Interpretation and Action} \\
\midrule
Tier 1A &
N57H, M66I &
Strong \claimlabel{S} &
High-priority confirmation regardless of subtype.
Bootstrap rank SD $<$ 0.5; structural and evolutionary concordance. \\
\midrule
Tier 1B &
K70N+N74K, Q67H+N74D &
Moderate \claimlabel{M} &
High-priority combination signals; confirm and document genetic context.
Q67H+N74D shows positive epistasis; K70N+N74K highest resistance without M66I. Bootstrap rank SD $\approx$ 0.8 (moderate stability). \\
\midrule
Tier 2 &
Q67H+K70R &
Hypothesis-generating \claimlabel{H} for mechanism; Moderate \claimlabel{M} for variability signal &
Flag for subtype/context-aware phenotyping. 78.5-fold range across studies indicates context sensitivity.
M66I+A105T, M66I+T107A: compensatory patterns tracked as fitness-restoring mutations. \\
\midrule
Tier 3 &
A105T/T107A compensatory pathways; subtype baseline variation &
Hypothesis-generating \claimlabel{H} &
Research priority. Not yet sufficient for clinical interpretation; inform prospective surveillance study design. \\
\bottomrule
\end{tabular}
\vspace{0.2cm}
\par\scriptsize\textit{Note:} Evidence strength evaluated across four dimensions: (1) phenotypic effect size and reproducibility; (2) bootstrap ranking stability; (3) structural plausibility; and (4) evolutionary constraint. Interpretive categories are intended as a provisional scaffold pending broader validation and complement expert-maintained HIV drug resistance mutation lists such as IAS-USA.
\end{table*}